% config/thesis-preamble.tex
% -----------------------------------------------------------------------------
% SHARED PREAMBLE
% -----------------------------------------------------------------------------
% Loaded by main.tex, proposal.tex, and skripsi-draft.tex right after
% \documentclass{template}. Every package, column type, and helper macro used by
% the manuscript lives here, so the three drivers differ only in which chapters
% they include.
%
% The driver declares its target before inputting this file:
%
%   \newcommand{\ThesisTarget}{proposal}   % proposal | draft | full
%   % config/thesis-preamble.tex
% -----------------------------------------------------------------------------
% SHARED PREAMBLE
% -----------------------------------------------------------------------------
% Loaded by main.tex, proposal.tex, and skripsi-draft.tex right after
% \documentclass{template}. Every package, column type, and helper macro used by
% the manuscript lives here, so the three drivers differ only in which chapters
% they include.
%
% The driver declares its target before inputting this file:
%
%   \newcommand{\ThesisTarget}{proposal}   % proposal | draft | full
%   % config/thesis-preamble.tex
% -----------------------------------------------------------------------------
% SHARED PREAMBLE
% -----------------------------------------------------------------------------
% Loaded by main.tex, proposal.tex, and skripsi-draft.tex right after
% \documentclass{template}. Every package, column type, and helper macro used by
% the manuscript lives here, so the three drivers differ only in which chapters
% they include.
%
% The driver declares its target before inputting this file:
%
%   \newcommand{\ThesisTarget}{proposal}   % proposal | draft | full
%   % config/thesis-preamble.tex
% -----------------------------------------------------------------------------
% SHARED PREAMBLE
% -----------------------------------------------------------------------------
% Loaded by main.tex, proposal.tex, and skripsi-draft.tex right after
% \documentclass{template}. Every package, column type, and helper macro used by
% the manuscript lives here, so the three drivers differ only in which chapters
% they include.
%
% The driver declares its target before inputting this file:
%
%   \newcommand{\ThesisTarget}{proposal}   % proposal | draft | full
%   \input{config/thesis-preamble}
%
%   proposal -> Bab I-III, no algorithm2e, no list of algorithms/appendices
%   draft    -> Bab I-IV, adds algorithm2e and both lists, adds the appendix
%   full     -> Bab I-V, adds all signature sheets, abstract, and preface
% -----------------------------------------------------------------------------

\providecommand{\ThesisTarget}{full}

\usepackage{ifthen}
\newif\ifThesisProposal
\newif\ifThesisDraft
\newif\ifThesisFull
\ifthenelse{\equal{\ThesisTarget}{proposal}}{\ThesisProposaltrue}{}
\ifthenelse{\equal{\ThesisTarget}{draft}}{\ThesisDrafttrue}{}
\ifthenelse{\equal{\ThesisTarget}{full}}{\ThesisFulltrue}{}
\ifThesisProposal\else\ifThesisDraft\else\ifThesisFull\else
  \ClassError{template}{Unknown \string\ThesisTarget: \ThesisTarget}%
    {Set \string\ThesisTarget\space to proposal, draft, or full.}%
\fi\fi\fi

% ── Mathematics
\usepackage{amsmath,amsfonts}
% Persamaan panjang boleh terpotong antarhalaman daripada meninggalkan
% halaman yang setengah kosong.
\allowdisplaybreaks[4]

% ── Graphics and floats
\usepackage{graphicx}
\usepackage{float}
\graphicspath{{figures/}}

% ── Tables
\usepackage{longtable}
\usepackage{tabularx}
\usepackage{xltabular}
\usepackage{booktabs}
% Kolom X berbobot dan rata kiri. Bobot seluruh kolom dalam satu tabel harus
% berjumlah sama dengan cacah kolomnya, misalnya 6 kolom -> total bobot 6.
% Rata kiri dipakai agar kolom sempit tidak diregangkan sampai berlubang.
\newcolumntype{L}[1]{>{\raggedright\arraybackslash\hsize=#1\hsize}X}
% Seluruh teks di dalam tabel dibuat rata kiri, bukan rata kanan-kiri.
% Kolom P menggantikan p{...} dan kolom Y menggantikan X.
\newcolumntype{P}[1]{>{\raggedright\arraybackslash}p{#1}}
\newcolumntype{Y}{>{\raggedright\arraybackslash}X}

% ── Pseudocode
% Naskah proposal tidak memuat algorithm2e: penyajian algoritma ditempatkan pada
% Bab IV, yang baru hadir pada naskah draf dan naskah skripsi penuh.
\ifThesisProposal\else
  \usepackage[ruled,vlined,linesnumbered]{algorithm2e}
\fi

% ── Source listings
% Paket listings tetap dimuat pada semua target karena config/thesis-language.tex
% mengubah \lstlistingname dan \lstlistlistingname pada awal dokumen.
\usepackage{listings}
\usepackage{xcolor}

% ── Cross-references and citations
\usepackage[hidelinks]{hyperref}
% Sitasi ganda pada satu kalimat diurutkan menaik dan ditulis dalam kurung
% terpisah, misalnya [4],[5] bukan [5,4]. Opsi nocompress mencegah rentang
% berurutan diringkas menjadi [4]-[6]. Paket cite dimuat setelah hyperref
% agar definisi \@citex miliknya tidak ditimpa kembali.
\usepackage[nocompress]{cite}
\renewcommand{\citepunct}{],[}

% ── Cover typeface and diagrams
\usepackage{helvet}
\usepackage{tikz}
\usetikzlibrary{calc}

% -----------------------------------------------------------------------------
% Configuration. The order matters: thesis-language.tex reads \ThesisLanguage
% and \ThesisDocumentType at load time, and \ThesisCoverTypeLabel reads the
% \ifThesis... flags declared above.
% -----------------------------------------------------------------------------
\input{config/thesis-config}
\input{config/thesis-language}

% ── Localized pseudocode keywords
\ifThesisProposal\else
  \SetKwInput{KwInput}{\Lang{Masukan}{Input}}
  \SetKwInput{KwOutput}{\Lang{Keluaran}{Output}}
  \SetKw{KwReturn}{\Lang{Kembalikan}{Return}}
\fi

\lstset{basicstyle=\ttfamily\small,breaklines=true,frame=single,columns=fullflexible,keywordstyle=\bfseries,captionpos=b}

% -----------------------------------------------------------------------------
% Writing-guide boxes and draft markers.
%
% \TemplateTip and \TemplateTodo explain what belongs in a section.
% \Placeholder marks prose that is still missing.
% \FigurePlaceholder{height}{description} stands in for an image that does not
% exist yet, inside a normal figure environment, so \caption and \label keep
% working and the figure keeps its number.
%
% All four are always defined, because the chapters call them unconditionally.
% \ThesisShowHelpers in config/thesis-config.tex decides whether they print.
% Remove every call before the final submission.
% -----------------------------------------------------------------------------
\newif\ifThesisHelpers
\ifthenelse{\equal{\ThesisShowHelpers}{on}}{\ThesisHelperstrue}{}
\ifthenelse{\equal{\ThesisShowHelpers}{auto}}{\ifThesisFull\ThesisHelperstrue\fi}{}

\ifThesisHelpers
  \newcommand{\TemplateTip}[1]{\par\noindent\fbox{\begin{minipage}{0.94\linewidth}\textbf{\Lang{Tips penulisan:}{Writing tip:}} #1\end{minipage}}\par\vspace{0.4em}}
  \newcommand{\TemplateTodo}[1]{\par\noindent\fbox{\begin{minipage}{0.94\linewidth}\textbf{\Lang{Isi bagian ini:}{Complete this section:}} #1\end{minipage}}\par\vspace{0.4em}}
  \newcommand{\Placeholder}[1]{\par\noindent\fbox{\begin{minipage}{0.94\linewidth}\textbf{[PLACEHOLDER]} #1\end{minipage}}\par\vspace{0.4em}}
  \newcommand{\FigurePlaceholder}[2]{%
    \setlength{\fboxsep}{0pt}%
    \fbox{\begin{minipage}[c][#1][c]{0.94\linewidth}%
    \centering\footnotesize\itshape\color{black!55}%
    \Lang{[GAMBAR BELUM TERSEDIA]}{[FIGURE NOT AVAILABLE YET]}\\[0.4em]#2%
    \end{minipage}}}
\else
  \newcommand{\TemplateTip}[1]{}
  \newcommand{\TemplateTodo}[1]{}
  \newcommand{\Placeholder}[1]{}
  \newcommand{\FigurePlaceholder}[2]{}
\fi

%
%   proposal -> Bab I-III, no algorithm2e, no list of algorithms/appendices
%   draft    -> Bab I-IV, adds algorithm2e and both lists, adds the appendix
%   full     -> Bab I-V, adds all signature sheets, abstract, and preface
% -----------------------------------------------------------------------------

\providecommand{\ThesisTarget}{full}

\usepackage{ifthen}
\newif\ifThesisProposal
\newif\ifThesisDraft
\newif\ifThesisFull
\ifthenelse{\equal{\ThesisTarget}{proposal}}{\ThesisProposaltrue}{}
\ifthenelse{\equal{\ThesisTarget}{draft}}{\ThesisDrafttrue}{}
\ifthenelse{\equal{\ThesisTarget}{full}}{\ThesisFulltrue}{}
\ifThesisProposal\else\ifThesisDraft\else\ifThesisFull\else
  \ClassError{template}{Unknown \string\ThesisTarget: \ThesisTarget}%
    {Set \string\ThesisTarget\space to proposal, draft, or full.}%
\fi\fi\fi

% ── Mathematics
\usepackage{amsmath,amsfonts}
% Persamaan panjang boleh terpotong antarhalaman daripada meninggalkan
% halaman yang setengah kosong.
\allowdisplaybreaks[4]

% ── Graphics and floats
\usepackage{graphicx}
\usepackage{float}
\graphicspath{{figures/}}

% ── Tables
\usepackage{longtable}
\usepackage{tabularx}
\usepackage{xltabular}
\usepackage{booktabs}
% Kolom X berbobot dan rata kiri. Bobot seluruh kolom dalam satu tabel harus
% berjumlah sama dengan cacah kolomnya, misalnya 6 kolom -> total bobot 6.
% Rata kiri dipakai agar kolom sempit tidak diregangkan sampai berlubang.
\newcolumntype{L}[1]{>{\raggedright\arraybackslash\hsize=#1\hsize}X}
% Seluruh teks di dalam tabel dibuat rata kiri, bukan rata kanan-kiri.
% Kolom P menggantikan p{...} dan kolom Y menggantikan X.
\newcolumntype{P}[1]{>{\raggedright\arraybackslash}p{#1}}
\newcolumntype{Y}{>{\raggedright\arraybackslash}X}

% ── Pseudocode
% Naskah proposal tidak memuat algorithm2e: penyajian algoritma ditempatkan pada
% Bab IV, yang baru hadir pada naskah draf dan naskah skripsi penuh.
\ifThesisProposal\else
  \usepackage[ruled,vlined,linesnumbered]{algorithm2e}
\fi

% ── Source listings
% Paket listings tetap dimuat pada semua target karena config/thesis-language.tex
% mengubah \lstlistingname dan \lstlistlistingname pada awal dokumen.
\usepackage{listings}
\usepackage{xcolor}

% ── Cross-references and citations
\usepackage[hidelinks]{hyperref}
% Sitasi ganda pada satu kalimat diurutkan menaik dan ditulis dalam kurung
% terpisah, misalnya [4],[5] bukan [5,4]. Opsi nocompress mencegah rentang
% berurutan diringkas menjadi [4]-[6]. Paket cite dimuat setelah hyperref
% agar definisi \@citex miliknya tidak ditimpa kembali.
\usepackage[nocompress]{cite}
\renewcommand{\citepunct}{],[}

% ── Cover typeface and diagrams
\usepackage{helvet}
\usepackage{tikz}
\usetikzlibrary{calc}

% -----------------------------------------------------------------------------
% Configuration. The order matters: thesis-language.tex reads \ThesisLanguage
% and \ThesisDocumentType at load time, and \ThesisCoverTypeLabel reads the
% \ifThesis... flags declared above.
% -----------------------------------------------------------------------------
% config/thesis-config.tex
% -----------------------------------------------------------------------------
% USER CONFIGURATION
% -----------------------------------------------------------------------------
% Change only this value to switch labels and selected content folders.
% The template loads files from content/<language>, chapters/<language>,
% and appendices/<language>.
% Valid values: indonesian, english
\newcommand{\ThesisLanguage}{indonesian}

% Change this value to choose the document type.
% Valid values: skripsi, tesis
% English labels: skripsi -> Undergraduate Thesis, tesis -> Thesis
\newcommand{\ThesisDocumentType}{skripsi}

% Visibility of the writing-guide boxes (\TemplateTip, \TemplateTodo) and the
% draft markers (\Placeholder, \FigurePlaceholder).
%   auto -> shown in main.tex, hidden in proposal.tex and skripsi-draft.tex
%   on   -> always shown
%   off  -> always hidden (use this for the final submitted manuscript)
% When hidden, \FigurePlaceholder collapses to nothing but the surrounding
% figure keeps its caption, number, and label.
\newcommand{\ThesisShowHelpers}{auto}

% Indonesian and English titles. Both are kept because Indonesian university
% documents commonly need both title forms even when the manuscript language changes.
\newcommand{\ThesisTitleID}{Judul Skripsi Bahasa Indonesia Ditulis dengan Huruf Kapital di Setiap Awal Kata}
\newcommand{\ThesisTitleEN}{English Thesis Title Written in Title Case}

\newcommand{\StudentName}{Nama Lengkap Mahasiswa}
\newcommand{\StudentID}{Nomor Pokok Mahasiswa}
\newcommand{\ProgramNameID}{Informatika}
\newcommand{\ProgramNameEN}{Informatics}

\newcommand{\AdvisorOne}{Nama Dosen Pembimbing I, Gelar}
\newcommand{\AdvisorOneID}{NIP/NPT. \makebox[4cm]{\dotfill}}
\newcommand{\AdvisorTwo}{Nama Dosen Pembimbing II, Gelar}
\newcommand{\AdvisorTwoID}{NIP/NPT. \makebox[4cm]{\dotfill}}
\newcommand{\ExaminerOne}{Nama Ketua Penguji, Gelar}
\newcommand{\ExaminerOneID}{NIP/NPT. \makebox[4cm]{\dotfill}}
\newcommand{\ExaminerTwo}{Nama Anggota Penguji, Gelar}
\newcommand{\ExaminerTwoID}{NIP/NPT. \makebox[4cm]{\dotfill}}

% Official metadata retained from the supplied template. Verify with the latest
% faculty/program guide before official submission.
\newcommand{\DeanName}{Nama Dekan Fakultas, Gelar}
\newcommand{\DeanID}{NIP./NPT. \makebox[4cm]{\dotfill}}
\newcommand{\HeadOfProgramName}{Nama Kepala Program Studi, Gelar}
\newcommand{\HeadOfProgramID}{NIP./NPT. \makebox[4cm]{\dotfill}}

\newcommand{\DefenseDateID}{Tanggal Bulan Tahun}
\newcommand{\DefenseDateEN}{Month Day, Year}
\newcommand{\SubmissionYear}{2026}
\newcommand{\KeywordsID}{kata kunci pertama, kata kunci kedua, kata kunci ketiga}
\newcommand{\KeywordsEN}{first keyword, second keyword, third keyword}

% config/thesis-language.tex
% -----------------------------------------------------------------------------
% LANGUAGE SWITCHING HELPERS
% -----------------------------------------------------------------------------
% This file converts \ThesisLanguage from config/thesis-config.tex into labels.
% Most students only need to edit \ThesisLanguage, not this file.

\RequirePackage{ifthen}

\newif\ifthesisenglish
\ifthenelse{\equal{\ThesisLanguage}{english}}{\thesisenglishtrue}{\thesisenglishfalse}

\newif\ifthesisdoctype
\ifthenelse{\equal{\ThesisDocumentType}{tesis}}{\thesisdoctypetrue}{\thesisdoctypefalse}

% Choose Indonesian text if \ThesisLanguage is indonesian, English text if english.
\newcommand{\Lang}[2]{\ifthesisenglish #2\else #1\fi}

% Derived metadata helpers.
\newcommand{\ThesisTitle}{\Lang{\ThesisTitleID}{\ThesisTitleEN}}
\newcommand{\ThesisTitleUpper}{\MakeUppercase{\ThesisTitle}}
\newcommand{\StudentNameUpper}{\MakeUppercase{\StudentName}}
\newcommand{\DefenseDate}{\Lang{\DefenseDateID}{\DefenseDateEN}}
\newcommand{\StudyProgramName}{\Lang{\ProgramNameID}{\ProgramNameEN}}
\newcommand{\StudyProgramPhrase}{\Lang{Program Studi \ProgramNameID}{\ProgramNameEN\ Study Program}}
\newcommand{\DocTypeNameID}{\ifthesisdoctype Tesis\else Skripsi\fi}
\newcommand{\DocTypeNameEN}{\ifthesisdoctype Thesis\else Undergraduate Thesis\fi}
\newcommand{\DocTypeNameLowerID}{\ifthesisdoctype tesis\else skripsi\fi}
\newcommand{\DocTypeNameLowerEN}{\ifthesisdoctype thesis\else undergraduate thesis\fi}
\newcommand{\DocTypeName}{\Lang{\DocTypeNameID}{\DocTypeNameEN}}
\newcommand{\DocTypeNameLower}{\Lang{\DocTypeNameLowerID}{\DocTypeNameLowerEN}}

% Document labels
\newcommand{\DocTypeLabel}{\DocTypeName}
\newcommand{\AdvisorBlockLabel}{\Lang{DOSEN PEMBIMBING}{ADVISORS}}
\newcommand{\MinistryLabel}{\Lang{KEMENTERIAN PENDIDIKAN TINGGI, SAINS, DAN TEKNOLOGI}{MINISTRY OF HIGHER EDUCATION, SCIENCE, AND TECHNOLOGY}}
\newcommand{\UniversityLabel}{\Lang{UNIVERSITAS PEMBANGUNAN NASIONAL VETERAN JAWA TIMUR}{UNIVERSITAS PEMBANGUNAN NASIONAL VETERAN JAWA TIMUR}}
\newcommand{\FacultyLabel}{\Lang{Fakultas Ilmu Komputer}{Faculty of Computer Science}}
\newcommand{\ProgramLabel}{\MakeUppercase{\StudyProgramPhrase}}
\newcommand{\ApprovalSheetLabel}{\Lang{LEMBAR PENGESAHAN}{APPROVAL SHEET}}
% LEMBAR PENGESAHAN (tim penguji) dan LEMBAR PERSETUJUAN (koordinator program
% studi) adalah dua halaman berbeda, sehingga judul Inggrisnya juga dibedakan.
\newcommand{\AcceptanceSheetLabel}{\Lang{LEMBAR PERSETUJUAN}{CONSENT SHEET}}
\newcommand{\PlagiarismStatementLabel}{\Lang{SURAT PERNYATAAN BEBAS PLAGIASI}{STATEMENT OF ORIGINALITY}}
\newcommand{\PrefaceLabel}{\Lang{KATA PENGANTAR}{ACKNOWLEDGEMENTS}}
\newcommand{\AbbreviationListLabel}{\Lang{DAFTAR SINGKATAN}{LIST OF ABBREVIATIONS}}
\newcommand{\GlossaryListLabel}{\Lang{GLOSARIUM}{GLOSSARY}}
\newcommand{\AppendixLabel}{\Lang{LAMPIRAN}{APPENDIX}}
\newcommand{\AppendixListLabel}{\Lang{DAFTAR LAMPIRAN}{LIST OF APPENDICES}}
\newcommand{\BibliographyLabel}{\Lang{DAFTAR PUSTAKA}{BIBLIOGRAPHY}}
% template.cls sudah menyediakan \ChapterPrefix sebagai cadangan agar berkas
% kelas itu tetap bisa dipakai sendiri, jadi di sini nilainya ditimpa.
\providecommand{\ChapterPrefix}{}
\renewcommand{\ChapterPrefix}{\Lang{BAB}{CHAPTER}}
\newcommand{\ContentsPrefix}{\Lang{BAB }{CHAPTER }}
\newcommand{\NpmLabel}{\Lang{NPM}{NPM}}
\newcommand{\ByLabel}{\Lang{Oleh}{By}}
\newcommand{\AdvisorOneRole}{\Lang{Pembimbing I}{Advisor I}}
\newcommand{\AdvisorTwoRole}{\Lang{Pembimbing II}{Advisor II}}
\newcommand{\ExaminerOneRole}{\Lang{Ketua Penguji}{Head Assessor}}
\newcommand{\ExaminerTwoRole}{\Lang{Anggota Penguji}{Assessor I}}
\newcommand{\AgreeLabel}{\Lang{Menyetujui,}{Approved by,}}
\newcommand{\KnowLabel}{\Lang{Mengetahui,}{Acknowledged by,}}
\newcommand{\DeanRoleLabel}{\Lang{Dekan Fakultas Ilmu Komputer}{Dean of the Faculty of Computer Science}}
\newcommand{\HeadProgramRoleLabel}{\Lang{Koordinator \StudyProgramPhrase}{Coordinator of \StudyProgramPhrase}}
\newcommand{\StudentNameLabel}{\Lang{Nama Mahasiswa}{Student Name}}
\newcommand{\ThesisTitleLabel}{\Lang{Judul \DocTypeName}{Thesis Title}}
\newcommand{\ProgramDegreeLabel}{\Lang{Sarjana (S1)}{Bachelor (S1)}}
\newcommand{\DegreeRequirement}{\Lang{\DocTypeName\ ini disusun sebagai salah satu syarat untuk memperoleh gelar Sarjana Komputer pada \StudyProgramPhrase\ Fakultas Ilmu Komputer Universitas Pembangunan Nasional Veteran Jawa Timur.}{This \DocTypeNameLower\ is prepared as one of the requirements for obtaining the Bachelor of Computer Science degree at the \StudyProgramPhrase, Faculty of Computer Science, Universitas Pembangunan Nasional Veteran Jawa Timur.}}


% Standard LaTeX names yang disesuaikan dengan dokumen asli
\AtBeginDocument{%
  \renewcommand{\figurename}{\Lang{Gambar}{Figure}}%
  \renewcommand{\tablename}{\Lang{Tabel}{Table}}%
  \renewcommand{\contentsname}{\Lang{DAFTAR ISI}{TABLE OF CONTENTS}}%
  \renewcommand{\listfigurename}{\Lang{DAFTAR GAMBAR}{LIST OF FIGURES}}%
  \renewcommand{\listtablename}{\Lang{DAFTAR TABEL}{LIST OF TABLES}}%
  \renewcommand{\bibname}{\BibliographyLabel}%
  \renewcommand{\lstlistingname}{\Lang{Kode Program}{Code Program}}%
  \renewcommand{\lstlistlistingname}{\Lang{DAFTAR KODE}{LIST OF CODE}}%
  % \providecommand dipasang lebih dahulu agar berkas ini tetap aman dimuat
  % dokumen yang tidak memuat algorithm2e, misalnya naskah proposal.
  \providecommand{\algorithmcfname}{}%
  \providecommand{\listalgorithmcfname}{}%
  \renewcommand{\algorithmcfname}{\Lang{Algoritma}{Algorithm}}%
  \renewcommand{\listalgorithmcfname}{\Lang{DAFTAR ALGORITMA}{LIST OF ALGORITHMS}}%
}

% Document-type wording printed on the cover. Each driver selects one through
% \ThesisTarget; see config/thesis-preamble.tex.
\newcommand{\ProposalTypeLabel}{\Lang{PRA-SKRIPSI}{RESEARCH PROPOSAL}}
\newcommand{\DraftTypeLabel}{\Lang{\DocTypeName\ (DRAF BAB I--IV)}{\DocTypeName\ (DRAFT CH. I--IV)}}
\newcommand{\ThesisCoverTypeLabel}{%
  \ifThesisProposal \ProposalTypeLabel
  \else\ifThesisDraft \DraftTypeLabel
  \else \DocTypeLabel
  \fi\fi}

% Reusable heading format for main chapters.
\newcommand{\ApplyMainChapterStyle}{%
  \titleformat{\chapter}[display]
    {\normalfont\normalsize\bfseries\centering}
    {\ChapterPrefix\ \thechapter}
    {0pt}
    {\MakeUppercase}%
  \titlespacing*{\chapter}{0pt}{0pt}{12pt}%
}


% ── Localized pseudocode keywords
\ifThesisProposal\else
  \SetKwInput{KwInput}{\Lang{Masukan}{Input}}
  \SetKwInput{KwOutput}{\Lang{Keluaran}{Output}}
  \SetKw{KwReturn}{\Lang{Kembalikan}{Return}}
\fi

\lstset{basicstyle=\ttfamily\small,breaklines=true,frame=single,columns=fullflexible,keywordstyle=\bfseries,captionpos=b}

% -----------------------------------------------------------------------------
% Writing-guide boxes and draft markers.
%
% \TemplateTip and \TemplateTodo explain what belongs in a section.
% \Placeholder marks prose that is still missing.
% \FigurePlaceholder{height}{description} stands in for an image that does not
% exist yet, inside a normal figure environment, so \caption and \label keep
% working and the figure keeps its number.
%
% All four are always defined, because the chapters call them unconditionally.
% \ThesisShowHelpers in config/thesis-config.tex decides whether they print.
% Remove every call before the final submission.
% -----------------------------------------------------------------------------
\newif\ifThesisHelpers
\ifthenelse{\equal{\ThesisShowHelpers}{on}}{\ThesisHelperstrue}{}
\ifthenelse{\equal{\ThesisShowHelpers}{auto}}{\ifThesisFull\ThesisHelperstrue\fi}{}

\ifThesisHelpers
  \newcommand{\TemplateTip}[1]{\par\noindent\fbox{\begin{minipage}{0.94\linewidth}\textbf{\Lang{Tips penulisan:}{Writing tip:}} #1\end{minipage}}\par\vspace{0.4em}}
  \newcommand{\TemplateTodo}[1]{\par\noindent\fbox{\begin{minipage}{0.94\linewidth}\textbf{\Lang{Isi bagian ini:}{Complete this section:}} #1\end{minipage}}\par\vspace{0.4em}}
  \newcommand{\Placeholder}[1]{\par\noindent\fbox{\begin{minipage}{0.94\linewidth}\textbf{[PLACEHOLDER]} #1\end{minipage}}\par\vspace{0.4em}}
  \newcommand{\FigurePlaceholder}[2]{%
    \setlength{\fboxsep}{0pt}%
    \fbox{\begin{minipage}[c][#1][c]{0.94\linewidth}%
    \centering\footnotesize\itshape\color{black!55}%
    \Lang{[GAMBAR BELUM TERSEDIA]}{[FIGURE NOT AVAILABLE YET]}\\[0.4em]#2%
    \end{minipage}}}
\else
  \newcommand{\TemplateTip}[1]{}
  \newcommand{\TemplateTodo}[1]{}
  \newcommand{\Placeholder}[1]{}
  \newcommand{\FigurePlaceholder}[2]{}
\fi

%
%   proposal -> Bab I-III, no algorithm2e, no list of algorithms/appendices
%   draft    -> Bab I-IV, adds algorithm2e and both lists, adds the appendix
%   full     -> Bab I-V, adds all signature sheets, abstract, and preface
% -----------------------------------------------------------------------------

\providecommand{\ThesisTarget}{full}

\usepackage{ifthen}
\newif\ifThesisProposal
\newif\ifThesisDraft
\newif\ifThesisFull
\ifthenelse{\equal{\ThesisTarget}{proposal}}{\ThesisProposaltrue}{}
\ifthenelse{\equal{\ThesisTarget}{draft}}{\ThesisDrafttrue}{}
\ifthenelse{\equal{\ThesisTarget}{full}}{\ThesisFulltrue}{}
\ifThesisProposal\else\ifThesisDraft\else\ifThesisFull\else
  \ClassError{template}{Unknown \string\ThesisTarget: \ThesisTarget}%
    {Set \string\ThesisTarget\space to proposal, draft, or full.}%
\fi\fi\fi

% ── Mathematics
\usepackage{amsmath,amsfonts}
% Persamaan panjang boleh terpotong antarhalaman daripada meninggalkan
% halaman yang setengah kosong.
\allowdisplaybreaks[4]

% ── Graphics and floats
\usepackage{graphicx}
\usepackage{float}
\graphicspath{{figures/}}

% ── Tables
\usepackage{longtable}
\usepackage{tabularx}
\usepackage{xltabular}
\usepackage{booktabs}
% Kolom X berbobot dan rata kiri. Bobot seluruh kolom dalam satu tabel harus
% berjumlah sama dengan cacah kolomnya, misalnya 6 kolom -> total bobot 6.
% Rata kiri dipakai agar kolom sempit tidak diregangkan sampai berlubang.
\newcolumntype{L}[1]{>{\raggedright\arraybackslash\hsize=#1\hsize}X}
% Seluruh teks di dalam tabel dibuat rata kiri, bukan rata kanan-kiri.
% Kolom P menggantikan p{...} dan kolom Y menggantikan X.
\newcolumntype{P}[1]{>{\raggedright\arraybackslash}p{#1}}
\newcolumntype{Y}{>{\raggedright\arraybackslash}X}

% ── Pseudocode
% Naskah proposal tidak memuat algorithm2e: penyajian algoritma ditempatkan pada
% Bab IV, yang baru hadir pada naskah draf dan naskah skripsi penuh.
\ifThesisProposal\else
  \usepackage[ruled,vlined,linesnumbered]{algorithm2e}
\fi

% ── Source listings
% Paket listings tetap dimuat pada semua target karena config/thesis-language.tex
% mengubah \lstlistingname dan \lstlistlistingname pada awal dokumen.
\usepackage{listings}
\usepackage{xcolor}

% ── Cross-references and citations
\usepackage[hidelinks]{hyperref}
% Sitasi ganda pada satu kalimat diurutkan menaik dan ditulis dalam kurung
% terpisah, misalnya [4],[5] bukan [5,4]. Opsi nocompress mencegah rentang
% berurutan diringkas menjadi [4]-[6]. Paket cite dimuat setelah hyperref
% agar definisi \@citex miliknya tidak ditimpa kembali.
\usepackage[nocompress]{cite}
\renewcommand{\citepunct}{],[}

% ── Cover typeface and diagrams
\usepackage{helvet}
\usepackage{tikz}
\usetikzlibrary{calc}

% -----------------------------------------------------------------------------
% Configuration. The order matters: thesis-language.tex reads \ThesisLanguage
% and \ThesisDocumentType at load time, and \ThesisCoverTypeLabel reads the
% \ifThesis... flags declared above.
% -----------------------------------------------------------------------------
% config/thesis-config.tex
% -----------------------------------------------------------------------------
% USER CONFIGURATION
% -----------------------------------------------------------------------------
% Change only this value to switch labels and selected content folders.
% The template loads files from content/<language>, chapters/<language>,
% and appendices/<language>.
% Valid values: indonesian, english
\newcommand{\ThesisLanguage}{indonesian}

% Change this value to choose the document type.
% Valid values: skripsi, tesis
% English labels: skripsi -> Undergraduate Thesis, tesis -> Thesis
\newcommand{\ThesisDocumentType}{skripsi}

% Visibility of the writing-guide boxes (\TemplateTip, \TemplateTodo) and the
% draft markers (\Placeholder, \FigurePlaceholder).
%   auto -> shown in main.tex, hidden in proposal.tex and skripsi-draft.tex
%   on   -> always shown
%   off  -> always hidden (use this for the final submitted manuscript)
% When hidden, \FigurePlaceholder collapses to nothing but the surrounding
% figure keeps its caption, number, and label.
\newcommand{\ThesisShowHelpers}{auto}

% Indonesian and English titles. Both are kept because Indonesian university
% documents commonly need both title forms even when the manuscript language changes.
\newcommand{\ThesisTitleID}{Judul Skripsi Bahasa Indonesia Ditulis dengan Huruf Kapital di Setiap Awal Kata}
\newcommand{\ThesisTitleEN}{English Thesis Title Written in Title Case}

\newcommand{\StudentName}{Nama Lengkap Mahasiswa}
\newcommand{\StudentID}{Nomor Pokok Mahasiswa}
\newcommand{\ProgramNameID}{Informatika}
\newcommand{\ProgramNameEN}{Informatics}

\newcommand{\AdvisorOne}{Nama Dosen Pembimbing I, Gelar}
\newcommand{\AdvisorOneID}{NIP/NPT. \makebox[4cm]{\dotfill}}
\newcommand{\AdvisorTwo}{Nama Dosen Pembimbing II, Gelar}
\newcommand{\AdvisorTwoID}{NIP/NPT. \makebox[4cm]{\dotfill}}
\newcommand{\ExaminerOne}{Nama Ketua Penguji, Gelar}
\newcommand{\ExaminerOneID}{NIP/NPT. \makebox[4cm]{\dotfill}}
\newcommand{\ExaminerTwo}{Nama Anggota Penguji, Gelar}
\newcommand{\ExaminerTwoID}{NIP/NPT. \makebox[4cm]{\dotfill}}

% Official metadata retained from the supplied template. Verify with the latest
% faculty/program guide before official submission.
\newcommand{\DeanName}{Nama Dekan Fakultas, Gelar}
\newcommand{\DeanID}{NIP./NPT. \makebox[4cm]{\dotfill}}
\newcommand{\HeadOfProgramName}{Nama Kepala Program Studi, Gelar}
\newcommand{\HeadOfProgramID}{NIP./NPT. \makebox[4cm]{\dotfill}}

\newcommand{\DefenseDateID}{Tanggal Bulan Tahun}
\newcommand{\DefenseDateEN}{Month Day, Year}
\newcommand{\SubmissionYear}{2026}
\newcommand{\KeywordsID}{kata kunci pertama, kata kunci kedua, kata kunci ketiga}
\newcommand{\KeywordsEN}{first keyword, second keyword, third keyword}

% config/thesis-language.tex
% -----------------------------------------------------------------------------
% LANGUAGE SWITCHING HELPERS
% -----------------------------------------------------------------------------
% This file converts \ThesisLanguage from config/thesis-config.tex into labels.
% Most students only need to edit \ThesisLanguage, not this file.

\RequirePackage{ifthen}

\newif\ifthesisenglish
\ifthenelse{\equal{\ThesisLanguage}{english}}{\thesisenglishtrue}{\thesisenglishfalse}

\newif\ifthesisdoctype
\ifthenelse{\equal{\ThesisDocumentType}{tesis}}{\thesisdoctypetrue}{\thesisdoctypefalse}

% Choose Indonesian text if \ThesisLanguage is indonesian, English text if english.
\newcommand{\Lang}[2]{\ifthesisenglish #2\else #1\fi}

% Derived metadata helpers.
\newcommand{\ThesisTitle}{\Lang{\ThesisTitleID}{\ThesisTitleEN}}
\newcommand{\ThesisTitleUpper}{\MakeUppercase{\ThesisTitle}}
\newcommand{\StudentNameUpper}{\MakeUppercase{\StudentName}}
\newcommand{\DefenseDate}{\Lang{\DefenseDateID}{\DefenseDateEN}}
\newcommand{\StudyProgramName}{\Lang{\ProgramNameID}{\ProgramNameEN}}
\newcommand{\StudyProgramPhrase}{\Lang{Program Studi \ProgramNameID}{\ProgramNameEN\ Study Program}}
\newcommand{\DocTypeNameID}{\ifthesisdoctype Tesis\else Skripsi\fi}
\newcommand{\DocTypeNameEN}{\ifthesisdoctype Thesis\else Undergraduate Thesis\fi}
\newcommand{\DocTypeNameLowerID}{\ifthesisdoctype tesis\else skripsi\fi}
\newcommand{\DocTypeNameLowerEN}{\ifthesisdoctype thesis\else undergraduate thesis\fi}
\newcommand{\DocTypeName}{\Lang{\DocTypeNameID}{\DocTypeNameEN}}
\newcommand{\DocTypeNameLower}{\Lang{\DocTypeNameLowerID}{\DocTypeNameLowerEN}}

% Document labels
\newcommand{\DocTypeLabel}{\DocTypeName}
\newcommand{\AdvisorBlockLabel}{\Lang{DOSEN PEMBIMBING}{ADVISORS}}
\newcommand{\MinistryLabel}{\Lang{KEMENTERIAN PENDIDIKAN TINGGI, SAINS, DAN TEKNOLOGI}{MINISTRY OF HIGHER EDUCATION, SCIENCE, AND TECHNOLOGY}}
\newcommand{\UniversityLabel}{\Lang{UNIVERSITAS PEMBANGUNAN NASIONAL VETERAN JAWA TIMUR}{UNIVERSITAS PEMBANGUNAN NASIONAL VETERAN JAWA TIMUR}}
\newcommand{\FacultyLabel}{\Lang{Fakultas Ilmu Komputer}{Faculty of Computer Science}}
\newcommand{\ProgramLabel}{\MakeUppercase{\StudyProgramPhrase}}
\newcommand{\ApprovalSheetLabel}{\Lang{LEMBAR PENGESAHAN}{APPROVAL SHEET}}
% LEMBAR PENGESAHAN (tim penguji) dan LEMBAR PERSETUJUAN (koordinator program
% studi) adalah dua halaman berbeda, sehingga judul Inggrisnya juga dibedakan.
\newcommand{\AcceptanceSheetLabel}{\Lang{LEMBAR PERSETUJUAN}{CONSENT SHEET}}
\newcommand{\PlagiarismStatementLabel}{\Lang{SURAT PERNYATAAN BEBAS PLAGIASI}{STATEMENT OF ORIGINALITY}}
\newcommand{\PrefaceLabel}{\Lang{KATA PENGANTAR}{ACKNOWLEDGEMENTS}}
\newcommand{\AbbreviationListLabel}{\Lang{DAFTAR SINGKATAN}{LIST OF ABBREVIATIONS}}
\newcommand{\GlossaryListLabel}{\Lang{GLOSARIUM}{GLOSSARY}}
\newcommand{\AppendixLabel}{\Lang{LAMPIRAN}{APPENDIX}}
\newcommand{\AppendixListLabel}{\Lang{DAFTAR LAMPIRAN}{LIST OF APPENDICES}}
\newcommand{\BibliographyLabel}{\Lang{DAFTAR PUSTAKA}{BIBLIOGRAPHY}}
% template.cls sudah menyediakan \ChapterPrefix sebagai cadangan agar berkas
% kelas itu tetap bisa dipakai sendiri, jadi di sini nilainya ditimpa.
\providecommand{\ChapterPrefix}{}
\renewcommand{\ChapterPrefix}{\Lang{BAB}{CHAPTER}}
\newcommand{\ContentsPrefix}{\Lang{BAB }{CHAPTER }}
\newcommand{\NpmLabel}{\Lang{NPM}{NPM}}
\newcommand{\ByLabel}{\Lang{Oleh}{By}}
\newcommand{\AdvisorOneRole}{\Lang{Pembimbing I}{Advisor I}}
\newcommand{\AdvisorTwoRole}{\Lang{Pembimbing II}{Advisor II}}
\newcommand{\ExaminerOneRole}{\Lang{Ketua Penguji}{Head Assessor}}
\newcommand{\ExaminerTwoRole}{\Lang{Anggota Penguji}{Assessor I}}
\newcommand{\AgreeLabel}{\Lang{Menyetujui,}{Approved by,}}
\newcommand{\KnowLabel}{\Lang{Mengetahui,}{Acknowledged by,}}
\newcommand{\DeanRoleLabel}{\Lang{Dekan Fakultas Ilmu Komputer}{Dean of the Faculty of Computer Science}}
\newcommand{\HeadProgramRoleLabel}{\Lang{Koordinator \StudyProgramPhrase}{Coordinator of \StudyProgramPhrase}}
\newcommand{\StudentNameLabel}{\Lang{Nama Mahasiswa}{Student Name}}
\newcommand{\ThesisTitleLabel}{\Lang{Judul \DocTypeName}{Thesis Title}}
\newcommand{\ProgramDegreeLabel}{\Lang{Sarjana (S1)}{Bachelor (S1)}}
\newcommand{\DegreeRequirement}{\Lang{\DocTypeName\ ini disusun sebagai salah satu syarat untuk memperoleh gelar Sarjana Komputer pada \StudyProgramPhrase\ Fakultas Ilmu Komputer Universitas Pembangunan Nasional Veteran Jawa Timur.}{This \DocTypeNameLower\ is prepared as one of the requirements for obtaining the Bachelor of Computer Science degree at the \StudyProgramPhrase, Faculty of Computer Science, Universitas Pembangunan Nasional Veteran Jawa Timur.}}


% Standard LaTeX names yang disesuaikan dengan dokumen asli
\AtBeginDocument{%
  \renewcommand{\figurename}{\Lang{Gambar}{Figure}}%
  \renewcommand{\tablename}{\Lang{Tabel}{Table}}%
  \renewcommand{\contentsname}{\Lang{DAFTAR ISI}{TABLE OF CONTENTS}}%
  \renewcommand{\listfigurename}{\Lang{DAFTAR GAMBAR}{LIST OF FIGURES}}%
  \renewcommand{\listtablename}{\Lang{DAFTAR TABEL}{LIST OF TABLES}}%
  \renewcommand{\bibname}{\BibliographyLabel}%
  \renewcommand{\lstlistingname}{\Lang{Kode Program}{Code Program}}%
  \renewcommand{\lstlistlistingname}{\Lang{DAFTAR KODE}{LIST OF CODE}}%
  % \providecommand dipasang lebih dahulu agar berkas ini tetap aman dimuat
  % dokumen yang tidak memuat algorithm2e, misalnya naskah proposal.
  \providecommand{\algorithmcfname}{}%
  \providecommand{\listalgorithmcfname}{}%
  \renewcommand{\algorithmcfname}{\Lang{Algoritma}{Algorithm}}%
  \renewcommand{\listalgorithmcfname}{\Lang{DAFTAR ALGORITMA}{LIST OF ALGORITHMS}}%
}

% Document-type wording printed on the cover. Each driver selects one through
% \ThesisTarget; see config/thesis-preamble.tex.
\newcommand{\ProposalTypeLabel}{\Lang{PRA-SKRIPSI}{RESEARCH PROPOSAL}}
\newcommand{\DraftTypeLabel}{\Lang{\DocTypeName\ (DRAF BAB I--IV)}{\DocTypeName\ (DRAFT CH. I--IV)}}
\newcommand{\ThesisCoverTypeLabel}{%
  \ifThesisProposal \ProposalTypeLabel
  \else\ifThesisDraft \DraftTypeLabel
  \else \DocTypeLabel
  \fi\fi}

% Reusable heading format for main chapters.
\newcommand{\ApplyMainChapterStyle}{%
  \titleformat{\chapter}[display]
    {\normalfont\normalsize\bfseries\centering}
    {\ChapterPrefix\ \thechapter}
    {0pt}
    {\MakeUppercase}%
  \titlespacing*{\chapter}{0pt}{0pt}{12pt}%
}


% ── Localized pseudocode keywords
\ifThesisProposal\else
  \SetKwInput{KwInput}{\Lang{Masukan}{Input}}
  \SetKwInput{KwOutput}{\Lang{Keluaran}{Output}}
  \SetKw{KwReturn}{\Lang{Kembalikan}{Return}}
\fi

\lstset{basicstyle=\ttfamily\small,breaklines=true,frame=single,columns=fullflexible,keywordstyle=\bfseries,captionpos=b}

% -----------------------------------------------------------------------------
% Writing-guide boxes and draft markers.
%
% \TemplateTip and \TemplateTodo explain what belongs in a section.
% \Placeholder marks prose that is still missing.
% \FigurePlaceholder{height}{description} stands in for an image that does not
% exist yet, inside a normal figure environment, so \caption and \label keep
% working and the figure keeps its number.
%
% All four are always defined, because the chapters call them unconditionally.
% \ThesisShowHelpers in config/thesis-config.tex decides whether they print.
% Remove every call before the final submission.
% -----------------------------------------------------------------------------
\newif\ifThesisHelpers
\ifthenelse{\equal{\ThesisShowHelpers}{on}}{\ThesisHelperstrue}{}
\ifthenelse{\equal{\ThesisShowHelpers}{auto}}{\ifThesisFull\ThesisHelperstrue\fi}{}

\ifThesisHelpers
  \newcommand{\TemplateTip}[1]{\par\noindent\fbox{\begin{minipage}{0.94\linewidth}\textbf{\Lang{Tips penulisan:}{Writing tip:}} #1\end{minipage}}\par\vspace{0.4em}}
  \newcommand{\TemplateTodo}[1]{\par\noindent\fbox{\begin{minipage}{0.94\linewidth}\textbf{\Lang{Isi bagian ini:}{Complete this section:}} #1\end{minipage}}\par\vspace{0.4em}}
  \newcommand{\Placeholder}[1]{\par\noindent\fbox{\begin{minipage}{0.94\linewidth}\textbf{[PLACEHOLDER]} #1\end{minipage}}\par\vspace{0.4em}}
  \newcommand{\FigurePlaceholder}[2]{%
    \setlength{\fboxsep}{0pt}%
    \fbox{\begin{minipage}[c][#1][c]{0.94\linewidth}%
    \centering\footnotesize\itshape\color{black!55}%
    \Lang{[GAMBAR BELUM TERSEDIA]}{[FIGURE NOT AVAILABLE YET]}\\[0.4em]#2%
    \end{minipage}}}
\else
  \newcommand{\TemplateTip}[1]{}
  \newcommand{\TemplateTodo}[1]{}
  \newcommand{\Placeholder}[1]{}
  \newcommand{\FigurePlaceholder}[2]{}
\fi

%
%   proposal -> Bab I-III, no algorithm2e, no list of algorithms/appendices
%   draft    -> Bab I-IV, adds algorithm2e and both lists, adds the appendix
%   full     -> Bab I-V, adds all signature sheets, abstract, and preface
% -----------------------------------------------------------------------------

\providecommand{\ThesisTarget}{full}

\usepackage{ifthen}
\newif\ifThesisProposal
\newif\ifThesisDraft
\newif\ifThesisFull
\ifthenelse{\equal{\ThesisTarget}{proposal}}{\ThesisProposaltrue}{}
\ifthenelse{\equal{\ThesisTarget}{draft}}{\ThesisDrafttrue}{}
\ifthenelse{\equal{\ThesisTarget}{full}}{\ThesisFulltrue}{}
\ifThesisProposal\else\ifThesisDraft\else\ifThesisFull\else
  \ClassError{template}{Unknown \string\ThesisTarget: \ThesisTarget}%
    {Set \string\ThesisTarget\space to proposal, draft, or full.}%
\fi\fi\fi

% ── Mathematics
\usepackage{amsmath,amsfonts}
% Persamaan panjang boleh terpotong antarhalaman daripada meninggalkan
% halaman yang setengah kosong.
\allowdisplaybreaks[4]

% ── Graphics and floats
\usepackage{graphicx}
\usepackage{float}
\graphicspath{{figures/}}

% ── Tables
\usepackage{longtable}
\usepackage{tabularx}
\usepackage{xltabular}
\usepackage{booktabs}
% Kolom X berbobot dan rata kiri. Bobot seluruh kolom dalam satu tabel harus
% berjumlah sama dengan cacah kolomnya, misalnya 6 kolom -> total bobot 6.
% Rata kiri dipakai agar kolom sempit tidak diregangkan sampai berlubang.
\newcolumntype{L}[1]{>{\raggedright\arraybackslash\hsize=#1\hsize}X}
% Seluruh teks di dalam tabel dibuat rata kiri, bukan rata kanan-kiri.
% Kolom P menggantikan p{...} dan kolom Y menggantikan X.
\newcolumntype{P}[1]{>{\raggedright\arraybackslash}p{#1}}
\newcolumntype{Y}{>{\raggedright\arraybackslash}X}

% ── Pseudocode
% Naskah proposal tidak memuat algorithm2e: penyajian algoritma ditempatkan pada
% Bab IV, yang baru hadir pada naskah draf dan naskah skripsi penuh.
\ifThesisProposal\else
  \usepackage[ruled,vlined,linesnumbered]{algorithm2e}
\fi

% ── Source listings
% Paket listings tetap dimuat pada semua target karena config/thesis-language.tex
% mengubah \lstlistingname dan \lstlistlistingname pada awal dokumen.
\usepackage{listings}
\usepackage{xcolor}

% ── Cross-references and citations
\usepackage[hidelinks]{hyperref}
% Sitasi ganda pada satu kalimat diurutkan menaik dan ditulis dalam kurung
% terpisah, misalnya [4],[5] bukan [5,4]. Opsi nocompress mencegah rentang
% berurutan diringkas menjadi [4]-[6]. Paket cite dimuat setelah hyperref
% agar definisi \@citex miliknya tidak ditimpa kembali.
\usepackage[nocompress]{cite}
\renewcommand{\citepunct}{],[}

% ── Cover typeface and diagrams
\usepackage{helvet}
\usepackage{tikz}
\usetikzlibrary{calc}

% -----------------------------------------------------------------------------
% Configuration. The order matters: thesis-language.tex reads \ThesisLanguage
% and \ThesisDocumentType at load time, and \ThesisCoverTypeLabel reads the
% \ifThesis... flags declared above.
% -----------------------------------------------------------------------------
% config/thesis-config.tex
% -----------------------------------------------------------------------------
% USER CONFIGURATION
% -----------------------------------------------------------------------------
% Change only this value to switch labels and selected content folders.
% The template loads files from content/<language>, chapters/<language>,
% and appendices/<language>.
% Valid values: indonesian, english
\newcommand{\ThesisLanguage}{indonesian}

% Change this value to choose the document type.
% Valid values: skripsi, tesis
% English labels: skripsi -> Undergraduate Thesis, tesis -> Thesis
\newcommand{\ThesisDocumentType}{skripsi}

% Visibility of the writing-guide boxes (\TemplateTip, \TemplateTodo) and the
% draft markers (\Placeholder, \FigurePlaceholder).
%   auto -> shown in main.tex, hidden in proposal.tex and skripsi-draft.tex
%   on   -> always shown
%   off  -> always hidden (use this for the final submitted manuscript)
% When hidden, \FigurePlaceholder collapses to nothing but the surrounding
% figure keeps its caption, number, and label.
\newcommand{\ThesisShowHelpers}{auto}

% Indonesian and English titles. Both are kept because Indonesian university
% documents commonly need both title forms even when the manuscript language changes.
\newcommand{\ThesisTitleID}{Judul Skripsi Bahasa Indonesia Ditulis dengan Huruf Kapital di Setiap Awal Kata}
\newcommand{\ThesisTitleEN}{English Thesis Title Written in Title Case}

\newcommand{\StudentName}{Nama Lengkap Mahasiswa}
\newcommand{\StudentID}{Nomor Pokok Mahasiswa}
\newcommand{\ProgramNameID}{Informatika}
\newcommand{\ProgramNameEN}{Informatics}

\newcommand{\AdvisorOne}{Nama Dosen Pembimbing I, Gelar}
\newcommand{\AdvisorOneID}{NIP/NPT. \makebox[4cm]{\dotfill}}
\newcommand{\AdvisorTwo}{Nama Dosen Pembimbing II, Gelar}
\newcommand{\AdvisorTwoID}{NIP/NPT. \makebox[4cm]{\dotfill}}
\newcommand{\ExaminerOne}{Nama Ketua Penguji, Gelar}
\newcommand{\ExaminerOneID}{NIP/NPT. \makebox[4cm]{\dotfill}}
\newcommand{\ExaminerTwo}{Nama Anggota Penguji, Gelar}
\newcommand{\ExaminerTwoID}{NIP/NPT. \makebox[4cm]{\dotfill}}

% Official metadata retained from the supplied template. Verify with the latest
% faculty/program guide before official submission.
\newcommand{\DeanName}{Nama Dekan Fakultas, Gelar}
\newcommand{\DeanID}{NIP./NPT. \makebox[4cm]{\dotfill}}
\newcommand{\HeadOfProgramName}{Nama Kepala Program Studi, Gelar}
\newcommand{\HeadOfProgramID}{NIP./NPT. \makebox[4cm]{\dotfill}}

\newcommand{\DefenseDateID}{Tanggal Bulan Tahun}
\newcommand{\DefenseDateEN}{Month Day, Year}
\newcommand{\SubmissionYear}{2026}
\newcommand{\KeywordsID}{kata kunci pertama, kata kunci kedua, kata kunci ketiga}
\newcommand{\KeywordsEN}{first keyword, second keyword, third keyword}

% config/thesis-language.tex
% -----------------------------------------------------------------------------
% LANGUAGE SWITCHING HELPERS
% -----------------------------------------------------------------------------
% This file converts \ThesisLanguage from config/thesis-config.tex into labels.
% Most students only need to edit \ThesisLanguage, not this file.

\RequirePackage{ifthen}

\newif\ifthesisenglish
\ifthenelse{\equal{\ThesisLanguage}{english}}{\thesisenglishtrue}{\thesisenglishfalse}

\newif\ifthesisdoctype
\ifthenelse{\equal{\ThesisDocumentType}{tesis}}{\thesisdoctypetrue}{\thesisdoctypefalse}

% Choose Indonesian text if \ThesisLanguage is indonesian, English text if english.
\newcommand{\Lang}[2]{\ifthesisenglish #2\else #1\fi}

% Derived metadata helpers.
\newcommand{\ThesisTitle}{\Lang{\ThesisTitleID}{\ThesisTitleEN}}
\newcommand{\ThesisTitleUpper}{\MakeUppercase{\ThesisTitle}}
\newcommand{\StudentNameUpper}{\MakeUppercase{\StudentName}}
\newcommand{\DefenseDate}{\Lang{\DefenseDateID}{\DefenseDateEN}}
\newcommand{\StudyProgramName}{\Lang{\ProgramNameID}{\ProgramNameEN}}
\newcommand{\StudyProgramPhrase}{\Lang{Program Studi \ProgramNameID}{\ProgramNameEN\ Study Program}}
\newcommand{\DocTypeNameID}{\ifthesisdoctype Tesis\else Skripsi\fi}
\newcommand{\DocTypeNameEN}{\ifthesisdoctype Thesis\else Undergraduate Thesis\fi}
\newcommand{\DocTypeNameLowerID}{\ifthesisdoctype tesis\else skripsi\fi}
\newcommand{\DocTypeNameLowerEN}{\ifthesisdoctype thesis\else undergraduate thesis\fi}
\newcommand{\DocTypeName}{\Lang{\DocTypeNameID}{\DocTypeNameEN}}
\newcommand{\DocTypeNameLower}{\Lang{\DocTypeNameLowerID}{\DocTypeNameLowerEN}}

% Document labels
\newcommand{\DocTypeLabel}{\DocTypeName}
\newcommand{\AdvisorBlockLabel}{\Lang{DOSEN PEMBIMBING}{ADVISORS}}
\newcommand{\MinistryLabel}{\Lang{KEMENTERIAN PENDIDIKAN TINGGI, SAINS, DAN TEKNOLOGI}{MINISTRY OF HIGHER EDUCATION, SCIENCE, AND TECHNOLOGY}}
\newcommand{\UniversityLabel}{\Lang{UNIVERSITAS PEMBANGUNAN NASIONAL VETERAN JAWA TIMUR}{UNIVERSITAS PEMBANGUNAN NASIONAL VETERAN JAWA TIMUR}}
\newcommand{\FacultyLabel}{\Lang{Fakultas Ilmu Komputer}{Faculty of Computer Science}}
\newcommand{\ProgramLabel}{\MakeUppercase{\StudyProgramPhrase}}
\newcommand{\ApprovalSheetLabel}{\Lang{LEMBAR PENGESAHAN}{APPROVAL SHEET}}
% LEMBAR PENGESAHAN (tim penguji) dan LEMBAR PERSETUJUAN (koordinator program
% studi) adalah dua halaman berbeda, sehingga judul Inggrisnya juga dibedakan.
\newcommand{\AcceptanceSheetLabel}{\Lang{LEMBAR PERSETUJUAN}{CONSENT SHEET}}
\newcommand{\PlagiarismStatementLabel}{\Lang{SURAT PERNYATAAN BEBAS PLAGIASI}{STATEMENT OF ORIGINALITY}}
\newcommand{\PrefaceLabel}{\Lang{KATA PENGANTAR}{ACKNOWLEDGEMENTS}}
\newcommand{\AbbreviationListLabel}{\Lang{DAFTAR SINGKATAN}{LIST OF ABBREVIATIONS}}
\newcommand{\GlossaryListLabel}{\Lang{GLOSARIUM}{GLOSSARY}}
\newcommand{\AppendixLabel}{\Lang{LAMPIRAN}{APPENDIX}}
\newcommand{\AppendixListLabel}{\Lang{DAFTAR LAMPIRAN}{LIST OF APPENDICES}}
\newcommand{\BibliographyLabel}{\Lang{DAFTAR PUSTAKA}{BIBLIOGRAPHY}}
% template.cls sudah menyediakan \ChapterPrefix sebagai cadangan agar berkas
% kelas itu tetap bisa dipakai sendiri, jadi di sini nilainya ditimpa.
\providecommand{\ChapterPrefix}{}
\renewcommand{\ChapterPrefix}{\Lang{BAB}{CHAPTER}}
\newcommand{\ContentsPrefix}{\Lang{BAB }{CHAPTER }}
\newcommand{\NpmLabel}{\Lang{NPM}{NPM}}
\newcommand{\ByLabel}{\Lang{Oleh}{By}}
\newcommand{\AdvisorOneRole}{\Lang{Pembimbing I}{Advisor I}}
\newcommand{\AdvisorTwoRole}{\Lang{Pembimbing II}{Advisor II}}
\newcommand{\ExaminerOneRole}{\Lang{Ketua Penguji}{Head Assessor}}
\newcommand{\ExaminerTwoRole}{\Lang{Anggota Penguji}{Assessor I}}
\newcommand{\AgreeLabel}{\Lang{Menyetujui,}{Approved by,}}
\newcommand{\KnowLabel}{\Lang{Mengetahui,}{Acknowledged by,}}
\newcommand{\DeanRoleLabel}{\Lang{Dekan Fakultas Ilmu Komputer}{Dean of the Faculty of Computer Science}}
\newcommand{\HeadProgramRoleLabel}{\Lang{Koordinator \StudyProgramPhrase}{Coordinator of \StudyProgramPhrase}}
\newcommand{\StudentNameLabel}{\Lang{Nama Mahasiswa}{Student Name}}
\newcommand{\ThesisTitleLabel}{\Lang{Judul \DocTypeName}{Thesis Title}}
\newcommand{\ProgramDegreeLabel}{\Lang{Sarjana (S1)}{Bachelor (S1)}}
\newcommand{\DegreeRequirement}{\Lang{\DocTypeName\ ini disusun sebagai salah satu syarat untuk memperoleh gelar Sarjana Komputer pada \StudyProgramPhrase\ Fakultas Ilmu Komputer Universitas Pembangunan Nasional Veteran Jawa Timur.}{This \DocTypeNameLower\ is prepared as one of the requirements for obtaining the Bachelor of Computer Science degree at the \StudyProgramPhrase, Faculty of Computer Science, Universitas Pembangunan Nasional Veteran Jawa Timur.}}


% Standard LaTeX names yang disesuaikan dengan dokumen asli
\AtBeginDocument{%
  \renewcommand{\figurename}{\Lang{Gambar}{Figure}}%
  \renewcommand{\tablename}{\Lang{Tabel}{Table}}%
  \renewcommand{\contentsname}{\Lang{DAFTAR ISI}{TABLE OF CONTENTS}}%
  \renewcommand{\listfigurename}{\Lang{DAFTAR GAMBAR}{LIST OF FIGURES}}%
  \renewcommand{\listtablename}{\Lang{DAFTAR TABEL}{LIST OF TABLES}}%
  \renewcommand{\bibname}{\BibliographyLabel}%
  \renewcommand{\lstlistingname}{\Lang{Kode Program}{Code Program}}%
  \renewcommand{\lstlistlistingname}{\Lang{DAFTAR KODE}{LIST OF CODE}}%
  % \providecommand dipasang lebih dahulu agar berkas ini tetap aman dimuat
  % dokumen yang tidak memuat algorithm2e, misalnya naskah proposal.
  \providecommand{\algorithmcfname}{}%
  \providecommand{\listalgorithmcfname}{}%
  \renewcommand{\algorithmcfname}{\Lang{Algoritma}{Algorithm}}%
  \renewcommand{\listalgorithmcfname}{\Lang{DAFTAR ALGORITMA}{LIST OF ALGORITHMS}}%
}

% Document-type wording printed on the cover. Each driver selects one through
% \ThesisTarget; see config/thesis-preamble.tex.
\newcommand{\ProposalTypeLabel}{\Lang{PRA-SKRIPSI}{RESEARCH PROPOSAL}}
\newcommand{\DraftTypeLabel}{\Lang{\DocTypeName\ (DRAF BAB I--IV)}{\DocTypeName\ (DRAFT CH. I--IV)}}
\newcommand{\ThesisCoverTypeLabel}{%
  \ifThesisProposal \ProposalTypeLabel
  \else\ifThesisDraft \DraftTypeLabel
  \else \DocTypeLabel
  \fi\fi}

% Reusable heading format for main chapters.
\newcommand{\ApplyMainChapterStyle}{%
  \titleformat{\chapter}[display]
    {\normalfont\normalsize\bfseries\centering}
    {\ChapterPrefix\ \thechapter}
    {0pt}
    {\MakeUppercase}%
  \titlespacing*{\chapter}{0pt}{0pt}{12pt}%
}


% ── Localized pseudocode keywords
\ifThesisProposal\else
  \SetKwInput{KwInput}{\Lang{Masukan}{Input}}
  \SetKwInput{KwOutput}{\Lang{Keluaran}{Output}}
  \SetKw{KwReturn}{\Lang{Kembalikan}{Return}}
\fi

\lstset{basicstyle=\ttfamily\small,breaklines=true,frame=single,columns=fullflexible,keywordstyle=\bfseries,captionpos=b}

% -----------------------------------------------------------------------------
% Writing-guide boxes and draft markers.
%
% \TemplateTip and \TemplateTodo explain what belongs in a section.
% \Placeholder marks prose that is still missing.
% \FigurePlaceholder{height}{description} stands in for an image that does not
% exist yet, inside a normal figure environment, so \caption and \label keep
% working and the figure keeps its number.
%
% All four are always defined, because the chapters call them unconditionally.
% \ThesisShowHelpers in config/thesis-config.tex decides whether they print.
% Remove every call before the final submission.
% -----------------------------------------------------------------------------
\newif\ifThesisHelpers
\ifthenelse{\equal{\ThesisShowHelpers}{on}}{\ThesisHelperstrue}{}
\ifthenelse{\equal{\ThesisShowHelpers}{auto}}{\ifThesisFull\ThesisHelperstrue\fi}{}

\ifThesisHelpers
  \newcommand{\TemplateTip}[1]{\par\noindent\fbox{\begin{minipage}{0.94\linewidth}\textbf{\Lang{Tips penulisan:}{Writing tip:}} #1\end{minipage}}\par\vspace{0.4em}}
  \newcommand{\TemplateTodo}[1]{\par\noindent\fbox{\begin{minipage}{0.94\linewidth}\textbf{\Lang{Isi bagian ini:}{Complete this section:}} #1\end{minipage}}\par\vspace{0.4em}}
  \newcommand{\Placeholder}[1]{\par\noindent\fbox{\begin{minipage}{0.94\linewidth}\textbf{[PLACEHOLDER]} #1\end{minipage}}\par\vspace{0.4em}}
  \newcommand{\FigurePlaceholder}[2]{%
    \setlength{\fboxsep}{0pt}%
    \fbox{\begin{minipage}[c][#1][c]{0.94\linewidth}%
    \centering\footnotesize\itshape\color{black!55}%
    \Lang{[GAMBAR BELUM TERSEDIA]}{[FIGURE NOT AVAILABLE YET]}\\[0.4em]#2%
    \end{minipage}}}
\else
  \newcommand{\TemplateTip}[1]{}
  \newcommand{\TemplateTodo}[1]{}
  \newcommand{\Placeholder}[1]{}
  \newcommand{\FigurePlaceholder}[2]{}
\fi
